\chapter{Conclusion and Future Work}
\label{chap:conclusion}

\section{Conclusion}

This thesis began with a practical tension. LoRA is simple and effective, but its rank and target
modules are fixed without reference to the downstream dataset. AdaLoRA makes allocation dynamic,
but adds schedule choices and relies on importance estimates during training. Star-LoRA explores a
middle path: use transparent dataset statistics to establish an allocation prior, then use
AdaLoRA for training.

The prototype measures dataset size, class imbalance, label entropy, token sparsity,
token-frequency skew, and embedding variance. These statistics control target rank, initial rank,
scaling, dropout, and adapter targets. The method also records a variance-normalized importance
score, although it does not yet use that score for allocation.

On Qwen2.5-0.5B and SST-2, Star-LoRA improves standard AdaLoRA most clearly at 1,000 and 5,000
examples. At 5,000 examples it reaches 91.55\% accuracy in the primary sweep, 4.97 percentage
points above AdaLoRA and slightly above LoRA. The selected rank increases with data volume, showing
that the policy responds as intended. However, fixed-rank LoRA remains stronger at several budgets
and is far faster. Star-LoRA's expanded adapter surface imposes a substantial cost.

The exploratory LLM-assisted policy does not improve accuracy. Claude produces a faster but weaker
layer layout, while Gemini requests fail and fall back to the heuristic. This result clarifies that
architectural suggestions without measured task feedback are insufficient for reliable allocation.

The appropriate conclusion is conditional. Dataset-aware configuration is a promising way to
improve AdaLoRA and reduce one-size-fits-all rank selection. The current evidence does not show
universal superiority over LoRA, and the method needs compute-aware design and controlled
ablations before stronger claims are justified.

\section{Future Work}

\subsection{Stability-Gated Allocation}

The most direct extension is to connect \cref{eq:stable-importance} to pruning. Components could be
eligible for removal only when their stable importance remains low for several intervals. Another
option is to blend PEFT's score with the normalized score:
\begin{equation}
  I^{\mathrm{combined}} =
  \lambda I^{\mathrm{PEFT}} +
  (1-\lambda)\widetilde{S}.
\end{equation}
This must be tested carefully because variance normalization can exaggerate parameters with very
small variance.

\subsection{Learned Dataset-to-Policy Mapping}

After collecting results from many datasets and ranks, a small regression or ranking model could
predict rank and module utility from profile features. Unlike the current hand-coded coefficients,
the controller could be trained against measured validation performance and constrained by cost.

\subsection{Cross-Task Evaluation}

Future experiments should include sentiment and topic classification, natural language inference,
imbalanced multi-class data, token classification, and generation. At least two model families and
multiple model scales are needed to evaluate whether the policy generalizes.

\subsection{Better Extreme-Low-Resource Behavior}

The policy could include a minimum-step rule. If the estimated training horizon is too short,
Star-LoRA could choose fixed LoRA instead of AdaLoRA. A meta-policy that selects among full
fine-tuning, LoRA, and AdaLoRA may be more useful than forcing one method into every regime.

\subsection{Compute-Aware Search}

Rank and module selection should be evaluated under equal parameter or runtime budgets. The policy
could estimate the parameter count of each candidate module and select the highest-value set under
a budget. Latency-aware objectives would make the method more relevant to deployment.

\subsection{Reproducibility Improvements}

Every future run should serialize all command-line arguments, package and CUDA versions, GPU model
and memory, git commit hash, dataset revision, trainable parameter count, profiling time, and peak
memory. These additions are inexpensive and would substantially strengthen subsequent analysis.
